% !TeX root = ../../Book4.tex
% 纲要标题：### 19.1 全导出函子
% 纲要标签：b4:sec:1801
% 纲要位置：计划纲要.md，第 3713 行。

\section{全导出函子}
\label{b4:sec:1801}

一个加性函子不一定保持拟同构，因而不能直接作用于导出态射。先选择能被它正确处理的复形模型，才能把导出群、自然变换与复合一起提升到导出范畴。

% 纲要内容（仅作写作计划，尚非教材正文）：
%
% * RF 和 LF 的构造及假设
% * 对象级导出函子与复形级导出的关系
% * 复合与自然变换
%

\subsection{\texorpdfstring{\(RF\) 与 \(LF\) 的构造条件}{RF 与 LF 的构造条件}}
\label{b4:subsec:1801:01}

对一个左正合函子逐次取导出，得到一列函子；但连接同态、自然变换及它们的复合仍需另外追踪。导出范畴允许保留产生这些群的整个复形。困难是原来的加性函子不一定保持拟同构，因而不能直接作用于分式表示。

\begin{definition}\label{b4:def:total-derived-functor}
设 \(\mathcal A,\mathcal B\) 是局部小交换范畴，\(F:\mathcal A\to\mathcal B\) 是加性函子，并固定导出范畴的集合大小约定。记 \(q_{\mathcal A}:K^+(\mathcal A)\to D^+(\mathcal A)\)、\(q_{\mathcal B}:K^+(\mathcal B)\to D^+(\mathcal B)\) 为局部化函子。三角函子
\(RF:D^+(\mathcal A)\to D^+(\mathcal B)\) 连同自然变换
\[
\eta:q_{\mathcal B}F\longrightarrow RFq_{\mathcal A}
\]
若满足下列泛性质，就称为 \(F\) 的\term[quan-you-daochu-hanzi]{全右导出函子}[total right derived functor]：对每个三角函子 \(T:D^+(\mathcal A)\to D^+(\mathcal B)\) 及自然变换 \(\xi:q_{\mathcal B}F\to Tq_{\mathcal A}\)，唯一存在自然变换 \(\theta:RF\to T\)，使 \(\xi=(\theta q_{\mathcal A})\eta\)。这里 \(F\) 在复形和同伦类上逐项作用。

对三角函子 \(LF:D^-(\mathcal A)\to D^-(\mathcal B)\)，将 \(q_{\mathcal A},q_{\mathcal B}\) 改为相应有上界范畴的局部化函子，并把自然变换方向改为
\(\epsilon:LFq_{\mathcal A}\to q_{\mathcal B}F\)。若对每个三角函子 \(T:D^-(\mathcal A)\to D^-(\mathcal B)\)，每个自然变换 \(Tq_{\mathcal A}\to q_{\mathcal B}F\) 都唯一经 \(T\to LF\) 分解，则称 \(LF\) 连同 \(\epsilon\) 为\term[quan-zuo-daochu-hanzi]{全左导出函子}[total left derived functor]。
\end{definition}
\symidx{\(RF,LF\)：全右导出函子与全左导出函子}

定义中的两个自然变换方向不同。内射替换给出 \(X\to I_X\)，因而右导出的比较来自 \(F(X)\to F(I_X)\)；投射替换给出 \(P_X\to X\)，因而左导出的比较来自 \(F(P_X)\to F(X)\)。加性函子未必保持这些拟同构，所以比较态射未必是同构。全导出函子修正的正是逐项 \(F\) 不能在导出范畴中直接使用的问题，不能要求修正前后的值总是相同。\cref{b4:prop:additive-right-derived-zero}将给出一个非零加性函子，其全右导出却为零。

\begin{theorem}\label{b4:thm:total-derived-functors}
设 \(\mathcal A,\mathcal B\) 是局部小交换范畴，\(F:\mathcal A\to\mathcal B\) 是加性函子，并在固定的集合大小约定下取导出范畴。
若 \(\mathcal A\) 有足够内射对象，则全右导出函子 \(RF:D^+(\mathcal A)\to D^+(\mathcal B)\) 存在，并对每个有下界上链复形 \(X\) 由
\[
RF(q_{\mathcal A}X)\cong q_{\mathcal B}F(I_X),\qquad X\xrightarrow{\simeq}I_X,
\]
计算，其中 \(I_X\) 是逐项内射的有下界复形。若 \(\mathcal A\) 有足够投射对象，则全左导出函子 \(LF:D^-(\mathcal A)\to D^-(\mathcal B)\) 存在，并对每个有上界上链复形 \(X\) 由
\[
LF(q_{\mathcal A}X)\cong q_{\mathcal B}F(P_X),\qquad P_X\xrightarrow{\simeq}X,
\]
计算，其中 \(P_X\) 是逐项投射的有上界复形。结果及其比较映射在导出范畴中不依赖替换的选择。
\end{theorem}
\begin{proof}
\cref{b4:thm:bounded-injective-model}给出的三角等价
\[
D^+(\mathcal A)\simeq K^+(\operatorname{Inj}\mathcal A)
\]
说明每个导出对象都有有下界的内射模型，而内射模型之间的导出态射可以按链同伦类计算。因此可先把导出对象换成内射模型，在模型上逐项作用 \(F\)，再进入 \(D^+(\mathcal B)\)。
记 \(\mathcal I=K^+(\operatorname{Inj}\mathcal A)\)，以
\(j:\mathcal I\to K^+(\mathcal A)\) 表示包含函子，并记此三角等价为
\(E=q_{\mathcal A}j:\mathcal I\to D^+(\mathcal A)\)。选择三角准逆
\(U:D^+(\mathcal A)\to\mathcal I\)，并将等价的自然同构取为伴随等价 \(U\dashv E\) 的单位与余单位：
\[
\rho:1_{D^+(\mathcal A)}\xrightarrow{\sim}EU,
\qquad
\kappa:UE\xrightarrow{\sim}1_{\mathcal I}.
\]
它们满足 \(E(\kappa_I)\rho_{E(I)}=1_{E(I)}\) 及
\(\kappa_{U(D)}U(\rho_D)=1_{U(D)}\)。定义
\(RF=q_{\mathcal B}FjU\)。加性函子保持有限双积、移位及锥的微分矩阵，所以逐项 \(F\) 是三角函子，所得复合也保持好三角。

对 \(X\in K^+(\mathcal A)\)，记 \(I_X=jU(q_{\mathcal A}X)\)。上一模型定理证明中的双射
\[
\operatorname{Hom}_{K^+(\mathcal A)}(X,I_X)
\xrightarrow{\sim}
\operatorname{Hom}_{D^+(\mathcal A)}(q_{\mathcal A}X,E U(q_{\mathcal A}X))
\]
将 \(\rho_{q_{\mathcal A}X}\) 唯一提升为同伦类 \(i_X:X\to I_X\)。因为其局部化是同构，\(i_X\) 的复形映射代表是拟同构。\(\rho\) 的自然性及上述双射的单射性，使 \(i\) 成为同伦范畴上的自然变换。这并不要求预先自然地选定严格链映射代表：更换代表只会改变一个链同伦，而加性 \(F\) 保持链同伦。因此
\(\eta_X=q_{\mathcal B}F(i_X)\) 不依赖代表，给出所需自然比较。

给定 \(T,\xi\)，先看泛分解等式怎样决定所需变换。第一条三角恒等式及提升唯一性给出 \(i_{jI}=j(\kappa_I^{-1})\)，第二条三角恒等式因而说明
\[
RF(\rho_D)=q_{\mathcal B}Fj(U(\rho_D))
=q_{\mathcal B}Fj(\kappa_{U(D)}^{-1})
=\eta_{jU(D)}.
\]
若 \(\theta:RF\to T\) 满足泛分解等式，沿 \(\rho_D:D\xrightarrow{\sim}EU(D)\) 的自然性便要求
\[
T(\rho_D)\theta_D
=\theta_{EU(D)}RF(\rho_D)
=\theta_{EU(D)}\eta_{jU(D)}
=\xi_{jU(D)}.
\]
由于 \(T(\rho_D)\) 可逆，\(\theta_D\) 只能取为
\[
\theta_D=(T(\rho_D))^{-1}\xi_{jU(D)}:
RF(D)\longrightarrow T(D).
\]
这里 \(\xi_{jU(D)}\) 的目标是 \(T(EU(D))\)；它与 \(T(\rho_D)\) 的逆可以复合。\(\xi\) 与 \(\rho\) 的自然性说明 \(\theta\) 自然，而
\[
\theta_{q_{\mathcal A}X}\eta_X
=(T(\rho_{q_{\mathcal A}X}))^{-1}
T(q_{\mathcal A}i_X)\xi_X=\xi_X.
\]
这证明存在性；前面的强制公式同时证明唯一性。

左导出采用 \(P_X\to X\) 的投射替换，所以比较方向是 \(LF\to F\)。给定的变换是 \(\xi:Tq_{\mathcal A}\to q_{\mathcal B}F\)，要找的是 \(\theta:T\to LF\)，使 \(\epsilon_X\theta_{q_{\mathcal A}X}=\xi_X\)；它的泛分解方向与右导出相反。记
\(\mathcal P=K^-(\operatorname{Proj}\mathcal A)\)，令
\(j^-:\mathcal P\to K^-(\mathcal A)\) 为包含，
\(E^-=q_{\mathcal A}j^-\)。由\cref{b4:thm:bounded-projective-model}选择三角准逆
\(V:D^-(\mathcal A)\to\mathcal P\)，以及伴随等价 \(E^-\dashv V\) 的单位、余单位
\[
\nu:1_{\mathcal P}\xrightarrow{\sim}VE^-,
\qquad \lambda:E^-V\xrightarrow{\sim}1_{D^-(\mathcal A)}.
\]
投射模型的 Hom 双射将 \(\lambda_{q_{\mathcal A}X}\) 唯一提升为自然的同伦类
\(p_X:j^-V(q_{\mathcal A}X)\to X\)。定义
\(LF=q_{\mathcal B}Fj^-V\)、\(\epsilon_X=q_{\mathcal B}F(p_X)\)。三角恒等式给出
\(LF(\lambda_D)=\epsilon_{j^-V(D)}\)。沿 \(\lambda_D:E^-V(D)\xrightarrow{\sim}D\) 的自然性和泛分解等式因而强制
\[
\theta_D T(\lambda_D)
=LF(\lambda_D)\theta_{E^-V(D)}
=\epsilon_{j^-V(D)}\theta_{E^-V(D)}
=\xi_{j^-V(D)}.
\]
所以所需变换在导出对象 \(D\) 上为
\[
\theta_D=\xi_{j^-V(D)}(T(\lambda_D))^{-1}:
T(D)\longrightarrow LF(D).
\]
自然性及 \(\epsilon_X\theta_{q_{\mathcal A}X}=\xi_X\) 由 \(\xi\) 与 \(\lambda\) 的自然性方块得到；前面的强制公式证明唯一性。

任意另一内射或投射替换都经模型的 Hom 双射得到与给定替换相容的唯一同伦比较；它在局部化后为同构，因模型等价全忠实而是同伦等价。逐项 \(F\) 保持同伦等价，故这些替换计算的对象典范同构于上述函子的值。泛性质也使不同模型选择所得函子连同比较变换唯一自然同构。
\end{proof}

这里构造存在只需加性，不要求左正合或右正合。正合性将用于把零次同调重新识别为 \(F\) 本身。以上泛性质及构造可对照 Weibel\cite[Definition 10.5.1, Existence Theorem 10.5.6]{Weibel1994HomologicalAlgebra}。

\subsection{对象级导出与复形级导出}
\label{b4:subsec:1801:02}

把对象 \(A\in\mathcal A\) 看作零次复形，替换 \(A\to I^\bullet\) 就是通常的内射消解。因此，当 \(F\) 左正合且 \(\mathcal A\) 有足够内射对象时，
\[
H^n(RF(A))\simeq R^nF(A)\quad(n\ge0),\qquad
H^n(RF(A))=0\quad(n<0).
\]
尤其 \(H^0(RF(A))=F(A)\)。右正合函子在有足够投射对象时给出
\(H^{-n}(LF(A))=L_nF(A)\)。负号来自本书始终使用升次微分，而通常的投射消解用非负链次数编号。

\ProseParagraphBreak
若 \(F\) 左正合、\(\mathcal A\) 有足够内射对象，而 \(X\) 为有下界复形，则 \(H^n(RF(X))\) 就是\cref{b4:def:hypercohomology}中的 \(\mathbb H^n(F,X)\)：两种计算选择同一内射复形后相同。\cref{b4:thm:hypercohomology-spectral-sequences}的第二套谱序列于是写成
\[
E_2^{p,q}=R^pF(H^q(X))
\Longrightarrow H^{p+q}(RF(X)).
\]
这里 \(p\ge0\)，\(q\) 有下界，每个总次数的滤过有限。即使知道所有第二页数据，仍须经过后续微分，并由终页数据得到目标的逐层商，不能一般地把它们直接相加。\(R^nF(A)\) 是单个对象的第 \(n\) 次导出，\(RF(X)\) 是整个复形的导出对象，\(H^n(RF(X))\) 才是该对象的第 \(n\) 次上同调。特别地，\(RF(X)\) 不能只由 \(R^nF(H^0X)\) 代替；即使 \(F\) 是恒等函子，\(RF(X)=X\) 仍保留每个 \(H^nX\)。

下文沿用\cref{b4:def:functor-acyclic}的“\(F\)-零调对象”：当 \(F\) 左正合时，它只要求 \(R^jF(Q)=0\) 对 \(j>0\) 成立；“无上同调对象”是同一性质的别称。它与整个复形的零调性是两个不同条件。

\begin{proposition}\label{b4:prop:total-derived-acyclic-complex}
设 \(F:\mathcal A\to\mathcal B\) 是局部小交换范畴间的加性左正合函子，\(\mathcal A\) 有足够内射对象。若 \(X\) 是 \(\mathcal A\) 中的有下界上链复形，且每项均满足
\(R^jF(X^n)=0\) 对所有 \(j>0\) 成立，则自然映射 \(F(X)\to RF(X)\) 是拟同构。若 \(F\) 为加性右正合函子，\(\mathcal A\) 有足够投射对象，\(X\) 为 \(\mathcal A\) 中有上界的上链复形，并且 \(L_jF(X^n)=0\) 对所有 \(n\in\mathbb Z\)、\(j>0\) 成立，则对偶的自然映射 \(LF(X)\to F(X)\) 是拟同构。
\end{proposition}
\begin{proof}
先设 \(E\) 是有下界的正合复形，每项为 \(F\)-零调对象。记 \(Z^n=\ker \mathrm{d}_E^n\)。从低于所有非零项的次数开始，\(Z^n=0\)。短正合列
\[
0\longrightarrow Z^n\longrightarrow E^n\longrightarrow Z^{n+1}\longrightarrow0
\]
给出导出长正合列。若 \(Z^n\) 已 \(F\)-零调，其中的片段
\[
F(E^n)\longrightarrow F(Z^{n+1})\longrightarrow R^1F(Z^n)=0
\]
使第一个映射满；而 \(E^n\) 的零调性给出
\[
R^iF(Z^{n+1})\cong R^{i+1}F(Z^n)=0\qquad(i\ge1),
\]
使下一层 \(Z^{n+1}\) 也 \(F\)-零调。左正合性再给出
\(0\to FZ^n\to FE^n\to FZ^{n+1}\to0\) 正合。由起始的 \(Z^n=0\) 逐次向上归纳，\(F(E)\) 便正合。

取 \(X\to I\) 的有下界内射替换，其锥正合、有下界，且每项是内射对象与某个 \(X^n\) 的有限直和，故逐项 \(F\)-零调。上段说明 \(F(\operatorname{Cone}(X\to I))\) 正合。加性保证此复形是 \(\operatorname{Cone}(FX\to FI)\)，锥判据得到拟同构。对偶情形从最高次数向下归纳。
\end{proof}

上一个证明说明了有下界假设的实际用途：它给余圈对象的归纳提供起点。逐项对 \(F\) 零调是对象性质，整个复形经 \(F\) 后仍正合是另一个性质；无界时不能仅凭前者推出后者。

\subsection{复合与自然变换}
\label{b4:subsec:1801:03}

若 \(\alpha:F\to G\) 是两个加性函子间的自然变换，在同一个内射替换上逐项取 \(\alpha_{I_X}\)，便得到
\(R\alpha:RF\to RG\)。比较态射使它自然，且
\(R(\beta\alpha)=(R\beta)(R\alpha)\)、\(R(1_F)=1_{RF}\)。左导出使用同一个投射替换，也有相同结论。因而自然变换不是在每个导出群上重新猜测，而是先在复形层面构造。

\begin{theorem}\label{b4:thm:total-derived-composition}
设 \(\mathcal A,\mathcal B,\mathcal C\) 是局部小交换范畴，\(\mathcal A,\mathcal B\) 有足够内射对象，\(F:\mathcal A\to\mathcal B\)、\(G:\mathcal B\to\mathcal C\) 为加性左正合函子。总有自然比较
\[
R(GF)\longrightarrow RG\circ RF
\quad\text{在}\;D^+(\mathcal A)\;\text{上}.
\]
若 \(F\) 把内射对象送到 \(G\)-零调对象，则这个比较是自然同构。
\end{theorem}
\begin{proof}
取 \(X\to I\)，再取 \(F(I)\to J\) 的有下界内射替换。\(R(GF)(X)\) 由 \(GF(I)\) 表示，\(RG(RF(X))\) 由 \(G(J)\) 表示，比较就是
\(G(F(I))\to G(J)\)：第一次导出使用 \(I\)，第二次导出还要把 \(F(I)\) 换成内射模型 \(J\)。同伦比较保证选择独立与自然性；等价地，它由全右导出的泛性质作用于两次比较变换的复合产生。虽然 \(F(I)\to J\) 为拟同构，加性左正合的 \(G\) 未必保持它，所以这个比较并不自动可逆。

若每个 \(F(I^n)\) 都 \(G\)-零调，则\cref{b4:prop:total-derived-acyclic-complex}应用于有下界复形 \(F(I)\)，说明上述比较是拟同构。这也给出两种复合方式之间的自然同构。
\end{proof}

若 \(F\) 把内射对象送到 \(G\)-零调对象，则对一般
\(X\in D^+(\mathcal A)\)，将\cref{b4:thm:hypercohomology-spectral-sequences}的第二套谱序列应用于 \(RF(X)\) 的有下界复形模型，得到自然的强收敛谱序列
\[
E_2^{p,q}=R^pG\bigl(H^q(RF(X))\bigr)
\Longrightarrow H^{p+q}\bigl(R(GF)(X)\bigr).
\]
这里 \(p\ge0\)，\(q\) 有下界，所以每个总次数的滤过仍有限。若特别取
\(X=A[0]\)，其中 \(A\in\mathcal A\)，则
\(H^q(RF(A[0]))=R^qF(A)\) 对 \(q\ge0\) 成立，负次数为零。于是恢复\cref{b4:thm:grothendieck-spectral-sequence}的第一象限谱序列
\[
E_2^{p,q}=R^pG(R^qF(A))
\Longrightarrow R^{p+q}(GF)(A).
\]
它提供目标的滤过及其逐层商，而不自动给出直和分解
\[
R^n(GF)(A)\cong\bigoplus_{p+q=n}R^pG(R^qF(A)).
\]
全导出复合的比较与零调性条件见 Weibel\cite[Composition Theorem 10.8.2, Corollary 10.8.3]{Weibel1994HomologicalAlgebra}。

\ProseParagraphBreak
对右正合函子及投射替换有对偶版本：若 \(F\) 把投射对象送到 \(G\)-零调对象，则 \(L(GF)\simeq LG\circ LF\)；一般比较的方向为 \(LG\circ LF\to L(GF)\)。如果不检查零调性，就没有理由让比较成为同构。

\begin{proposition}\label{b4:prop:additive-right-derived-zero}
在交换群范畴 \(\mathbf{Ab}\) 上取加性函子
\[
F(M)=M\otimes_{\mathbb Z}\mathbb Z/2\mathbb Z.
\]
其全右导出函子 \(RF:D^+(\mathbf{Ab})\to D^+(\mathbf{Ab})\) 为零函子，但 \(F(\mathbb Z)=\mathbb Z/2\mathbb Z\ne0\)。因此在零次复形 \(\mathbb Z[0]\) 上，比较态射 \(F(\mathbb Z)[0]\to RF(\mathbb Z[0])\) 不是同构。仅有加性不能保证对每个交换群 \(M\) 都有 \(H^0(RF(M[0]))\cong F(M)\)。
\end{proposition}
\begin{proof}
设 \(I\) 是内射交换群。对任意 \(x\in I\)，将同态 \(2\mathbb Z\to I\)、\(2\mapsto x\) 沿包含 \(2\mathbb Z\hookrightarrow\mathbb Z\) 延拓为 \(f:\mathbb Z\to I\)。令 \(y=f(1)\)，便有 \(2y=f(2)=x\)。故乘 \(2\) 在 \(I\) 上满射，\(I/2I=0\)，从而
\[
F(I)\cong I/2I=0.
\]
对任意有下界复形 \(X\)，取有下界的逐项内射替换 \(X\to I_X\)。\(F(I_X)\) 逐项为零，所以\cref{b4:thm:total-derived-functors}给出的 \(RF\) 是零函子。另一方面，\(F(\mathbb Z)=\mathbb Z/2\mathbb Z\) 非零，其零次复形的 \(H^0\) 非零，不能同构于零导出对象 \(RF(\mathbb Z[0])\)。该函子为右正合函子，但它把单射 \(\mathbb Z\xrightarrow{2}\mathbb Z\) 送为 \(\mathbb Z/2\mathbb Z\) 上的零映射，故不是左正合函子。因此它的全右导出不受前面左正合情形的零次识别约束。
\end{proof}
